\section{Gradient-based methods: Lucas--Kanade, Horn--Schunck}\label{sec:T31}
% verbatim from SVP/work/endsem/atoms/notes/T31.tex lines 111-113 (2026-11-06): a \flashcard whose back holds
% maths environments with '\\v' row breaks. web.py split the back off the command (render failure
% "\begin{minipage} ended by \end{preview}") and read '\\v' as the macro \v (math fallback).
\flashcard{Lucas--Kanade normal equations?}{$\begin{bmatrix}\sum I_x^2&\sum I_xI_y\\\sum I_xI_y&\sum I_y^2\end{bmatrix}\begin{bmatrix}u\\v\end{bmatrix}=-\begin{bmatrix}\sum I_xI_t\\\sum I_yI_t\end{bmatrix}$ over the window.}
\flashcard{When is the Lucas--Kanade system well conditioned?}{When both eigenvalues of the structure tensor $G$ are large: a corner or textured patch. An edge gives one near-zero eigenvalue.}
\flashcard{Horn--Schunck energy?}{$\iint(I_xu+I_yv+I_t)^2+\alpha^2(\lVert\nabla u\rVert^2+\lVert\nabla v\rVert^2)$: data term plus smoothness weighted by $\alpha^2$.}

% T31.tex lines 66-69: an inline list of tuples is wider than a 360 px phone line (page scrolled sideways)
\paragraph{Worked example.} A window of 4 pixels has $(I_x,I_y,I_t)=(1,0,-2),(0,1,-1),(1,1,-3),(2,1,-5)$.
$\sum I_x^2=6$, $\sum I_xI_y=3$, $\sum I_y^2=3$, $\sum I_xI_t=-15$, $\sum I_yI_t=-9$. So
$\left[\begin{smallmatrix}6&3\\3&3\end{smallmatrix}\right]\left[\begin{smallmatrix}u\\v\end{smallmatrix}\right]=\left[\begin{smallmatrix}15\\9\end{smallmatrix}\right]$,
$\det G=18-9=9$, $u=(3\cdot15-3\cdot9)/9=2$, $v=(6\cdot9-3\cdot15)/9=1$.

Order 1: $a_1=R(1)/R(0)$. Order 2: $\begin{bmatrix}R(0)&R(1)\\R(1)&R(0)\end{bmatrix}\begin{bmatrix}a_1\\a_2\end{bmatrix}
=\begin{bmatrix}R(1)\\R(2)\end{bmatrix}$. Example $R(0)=1$, $R(1)=0.8$, $R(2)=0.5$: determinant $1-0.64=0.36$,
$a_1=(0.8\cdot1-0.8\cdot0.5)/0.36=$ \answer{1.111}, $a_2=(1\cdot0.5-0.8^2)/0.36=$ \answer{$-0.389$}.

\paragraph{Worked example (Expected).} $H=\begin{psmallmatrix}1.2&0.1&30\\0.05&1.1&10\\0.001&0&1\end{psmallmatrix}$, point
$(100,50)$: $u=120+5+30=155$, $v=5+55+10=70$, $w=0.1+0+1=1.1$, so \answer{$(x',y')=(140.91,\,63.64)$}. Forgetting the
division by $w$ gives the wrong $(155,70)$.

\textbf{Practice numericals (notes p9), solved.}\par\nobreak
\begin{tabularx}{\linewidth}{@{}L{50mm}cccY@{}}\toprule
\thead{Signal} & \thead{$N$} & \thead{$L$} & \thead{$H$} & \thead{frames}\\\midrule
speech & 8000 & 200 & 80 & 98\\\bottomrule
\end{tabularx}
